\chapter{数值方法、计算细节与可复现实验设计}\label{ch:numerics}
\section{先写方程合同，再选求解器}
模拟论文中的“two-fluid code”不足以复现实验。必须记录：是否准中性、是否保留位移电流、电子惯性的完整形式、压力是标量还是张量、热流闭合、碰撞／超电阻／滤波项，以及它们进入哪条方程。边界和归一化同样属于模型定义。特别是 $m_i/m_e$、$c/\VA$、$k_0d_e$、$\eta$、$\Delta x/d_e$ 不能省略。

\begin{longtable}{P{30mm}P{54mm}P{64mm}}
\toprule 路线 & 数值核心 & 易混淆的误差\\\midrule
F3D类差分双流体 & 在规则网格上演化流体和磁场，加入Hall与选定电子惯性 & 人工扩散可能选定电子层长度；周期边界限制持续流出\\
有限体积Hall MHD & 守恒通量加Hall电场，配CT或散度控制 & Hall项含空间二阶导，普通理想MHD波速不足以定时间步\\
隐式扩展MHD & Newton--Krylov或相关非线性迭代、物理预条件 & 时间步可稳定却不能正确解析波相位或增长率\\
五矩／十矩Maxwell & 双曲输运与电磁源项，常配分裂、IMEX或局地隐式源项 & 光波、等离子频率、回旋频率刚性；压力正定性\\
谱／伪谱约化模型 & FFT求导、Poisson及Helmholtz反演、去混叠 & 周期性和光滑性限制；小尺度能量堆积及滤波耗散\\
有限元扩展MHD & 适合曲面和托卡马克几何，半隐式／隐式时间推进 & 网格各向异性、几何度量与闭合参数共同影响撕裂\\\bottomrule
\end{longtable}
方法入口见 \citet{Cassak2008,Toth2008,Arnold2008,Hakim2008,Wang2015}。不能将某方法论文证明的单波稳定性，等同于重联非线性收敛。

\section{Hall感应方程如何离散}
\subsection{交错电磁网格与约束输运}
二维示例中，将 $B_x$ 放在垂直 $x$ 的面、$B_y$ 放在垂直 $y$ 的面，$E_z$ 放在边。Faraday更新为
\begin{align}
B_{x,i+1/2,j}^{n+1}&=B_{x,i+1/2,j}^n-\frac{\Delta t}{\Delta y}(E_{z,i+1/2,j+1/2}-E_{z,i+1/2,j-1/2}),\\
B_{y,i,j+1/2}^{n+1}&=B_{y,i,j+1/2}^n+\frac{\Delta t}{\Delta x}(E_{z,i+1/2,j+1/2}-E_{z,i-1/2,j+1/2}).
\end{align}
对这些面磁场计算离散散度，相邻边电场成对抵消，所以散度在舍入误差内保持。这是CT的代数核心。Hall电场 $\E_H=\J\times\B/(en)$ 必须先在一致位置计算电流、磁场和密度，再插值到边；低密度区 $1/n$ 会放大误差。密度下限应作为模型干预记录并扫描，不能默认为无影响。

\subsection{Hall项不是普通耗散通量}
理想MHD常用近似Riemann求解器给出对流通量；Hall项包含 $\J=\grad\times\B/\mu_0$，在感应方程中形成二阶、色散性的空间算子。若把它简单当作普通扩散，可能阻尼掉真实whistler分支。正确的验收包含频率、群速度、极化和幅度误差，而不仅看计算没有发散。

对规则中心差分，Hall半离散谱在无耗散理想情况下近于虚轴；时间推进稳定域必须覆盖该谱。不同Runge--Kutta、预测校正和子循环方案的CFL常数不同，不能只抄一个通用 $C=0.5$。

\subsection{时间步预算}
显式方案可按最紧的约束估计
\begin{equation}\Delta t\lesssim\min\left[C_f{\Delta x\over |u|+c_f},\ C_H{\Delta x^2\over \VA d_i},\ C_\eta{\Delta x^2\over D_\eta},\ C_\nu{\Delta x^2\over\nu},\ C_c{\Delta x\over c},\ C_p\omega_{pe}^{-1}\right].
\end{equation}
最后两项只对显式全Maxwell／相关显式源更新适用。十矩磁旋转源还可能受 $\Omega_{ce}^{-1}$ 限制。隐式或解析源积分可放宽稳定性限制，但要保持电子响应、色散和能量交换的精度。

\section{电子惯性、隐式求解与守恒}
常密度EMHD可以演化 $\B_e$，再解
\begin{equation}(1-d_e^2\nabla^2)\B=\B_e.
\end{equation}
周期Fourier空间中 $\hat\B=\hat\B_e/(1+d_e^2k^2)$，是逐模代数反演；非周期需给出一致的边界条件。变密度和一般扩展MHD产生变系数算子，不能照搬常系数逆。

全隐式时间离散可写为 $F(U^{n+1})=0$。Newton步解 $\mathcal J\delta U=-F$，Krylov方法只需Jacobian向量积；物理预条件应近似Hall的whistler刚性和流体耦合。应报告每步非线性残差、线性迭代数、失败步回退以及能量残差。稳定大步长可能过度平滑线性撕裂，必须对时间步进行独立收敛检验。\citep{Arnold2008,Toth2008}

\section{五矩与十矩实现：源分裂、压力正定性和闭合}
双曲通量部分可用有限体积或DG。源部分包括 $q_sn_s\E$ 和磁旋转，且电场受所有粒子电流反馈；分开显式处理可能引入非物理等离子振荡。局地耦合隐式更新或一致的IMEX可缓解刚性。若采用Strang分裂，源半步、输运整步、源半步并不自动保证在强刚性极限保持二阶，需专门验证。

十矩压力正定性要求对任意非零 $\vct a$，$a_iP_{ij}a_j>0$。只检查三个对角分量为正不够，因为过大非对角分量会产生负特征值。每步应统计最小特征值、修正次数与修正能量。大量触发压力修正的算例不能仅凭平滑图像宣布可信。

局地闭合有精确松弛子步：$\mathsf P-p\mathsf I\mapsto e^{-\nu\Delta t}(\mathsf P-p\mathsf I)$，迹不变。温度梯度闭合需边界热通量；非局地闭合需额外方向算子。不同闭合比较必须在相同网格、初值、质量比、热参数、边界条件下进行，否则不能把差异归因于闭合本身。\citep{Wang2015,Ng2020,Ng2017,Closure2026}

\section{Harris电流片与GEM：把初始条件真正写全}
经典单Harris片可写为
\begin{equation}B_x=B_0\tanh(y/a),\quad B_z=B_g,\quad
n=n_b+n_h\operatorname{sech}^2(y/a),\quad
n_h(T_i+T_e)=B_0^2/(2\mu_0).
\end{equation}
因此 $p_i+p_e+B_x^2/(2\mu_0)$ 为常数。由Ampère定律，
\begin{equation}J_z=-{B_0\over\mu_0a}\operatorname{sech}^2(y/a).
\end{equation}
若按各物种的垂直力平衡分配电流且 $E_y=0$，取
\begin{equation}u_{iz}={2T_i\over eB_0a}{n_h\operatorname{sech}^2(y/a)\over n}\,(-1),\qquad
u_{ez}={2T_e\over eB_0a}{n_h\operatorname{sech}^2(y/a)\over n}.
\end{equation}
这里离子为负 $z$ 漂移、电子为正 $z$ 漂移，从而 $en(u_{iz}-u_{ez})=J_z$。在严格动理学Harris分布中，可把背景粒子与漂移电流片粒子分别构造；在流体初值中则要用总密度对应的平均速度。

GEM比较的经典归一化设置包括 $a=0.5d_i$、$n_b/n_h=0.2$、$T_i/T_e=5$，常用 $L_x=25.6d_i,L_y=12.8d_i$、质量比25、$x$周期及 $y$导体边界。初始扰动由磁矢势给出，例如
\begin{equation}\delta A_z=A_0\cos(2\pi x/L_x)\cos(\pi y/L_y),\quad \delta\B=\grad\times(\delta A_z\hat z),
\end{equation}
常见挑战振幅为 $A_0=0.1B_0d_i$。不同参赛模型的惯性、耗散和网格并不全部相同，实际复现须按具体论文核对。\citep{Birn2001}

\note{“平衡”必须检查各物种方程残差，而不只是总压强平衡。把全部电流分给电子，却仍宣称 $E_y=0$ 且离子存在压强梯度，会启动额外瞬态。强种子GEM适合比较已启动的非线性阶段，不适合单独证明无扰动系统何时自发触发。}

\section{建议的计算参数与资源估算}
下面是本报告建议的教学与研究起点，不是宣称已运行的模拟。先用五矩或明确耗散的Hall模型复现单片，再升级十矩。取 $L_x\times L_y=25.6\times12.8d_i^2$、$m_i/m_e=100$，为 $d_e$ 放4、8、16格时，网格分别为 $1024\times512$、$2048\times1024$、$4096\times2048$。按每格26个双精度变量，单份状态分别约104 MiB、416 MiB、1.63 GiB；包含重构、Runge--Kutta暂存和通信后常需数倍内存。完整三维会迅速超过桌面计算预算。

选择 $c/\VA$ 时同时检查 $v_{te}/c\ll1$ 和 $\lambda_{De}$ 的模型处理。全Maxwell双流体不是PIC，不必把所有PIC噪声准则照搬，但未解析的电荷分离波仍可能损害显式流体稳定性或精度。准中性模型则从方程层面删去了这些波，两个方案不能仅靠减小 $c$ 相互替代。

\section{重联诊断：具体量怎样计算}
\subsection{磁通与电场}
二维磁通函数 $A_z$ 的X点与相邻O点差定义重联磁通 $\Psi=A_z(O)-A_z(X)$。若临界点随时间移动，$\grad A_z=0$ 使轨迹导数中的平移项消失，
\begin{equation}\dot\Psi=-E_z(O)+E_z(X).
\end{equation}
并不是所有边界下 $\dot\Psi=E_z(X)$；只有参考O点电场可忽略等条件下才成立。多个X/O点出现后需跟踪拓扑、区分净磁通与反复重联的总通量。

\subsection{欧姆定律逐项闭合}
保存 $\E$、$\vel_e$、$\vel_i$、$n$、$\J$、全部压力分量及时间导数。计算
\begin{equation}\vct\epsilon_{\rm Ohm}=\E+\vel_i\times\B-{\J\times\B\over en}+{\grad\cdot\mathsf P_e\over en}+{m_e\over e}D_e\vel_e-\eta\J.
\end{equation}
这应接近离散截断误差，而不是接近某个物理项的量级。若 $E_z$ 的主要来源恰好是未记账残差，所谓无碰撞重联机制没有被证明。时间差分的采样间隔要足够细，否则会把电子惯性误差当成压力支撑。

\subsection{能量、谱与三维指标}
记录边界Poynting通量、两种粒子体动能与内能、$-\mathsf P_s:\grad\vel_s$、电阻热和人工耗散。三维可补充场线积分 $\int E_\parallel\dd l$、连接映射、磁通绳寿命、场线随机性和各向异性谱；有磁零点或开放场线时积分域及端点要明示。三维全局重联不能只用某一切片的 $E_z(X)$。

\section{最小验证矩阵}
\begin{enumerate}
\item 均匀平衡与磁旋转：验证源项、压力迹及Gauss／磁散度约束。
\item whistler色散：测量多个 $kd_i$ 的频率和阻尼，包含近电子尺度。
\item 线性撕裂：对扰动幅度指数拟合 $\gamma$，同时改变网格和时间步。
\item GEM：比较磁通曲线、电子／离子流出、Hall场及欧姆定律分解。
\item 岛合并：改变系统尺寸而保持局地模型参数，测量峰值与时间平均率。
\item 三维：固定横向分辨率，独立改变第三维长度及模式数，区分二维限制与真实不稳定性。
\end{enumerate}
验证表中的“收敛”至少指网格、时间步、模型参数三类依赖被区分。只在两张看起来相似的 $J_z$ 图之间比较，无法判定电子层或高阶压力已收敛。

\section{已核读论文实际采用了什么：四条数值路线}
\subsection{F3D：常系数惯性反演与人工四阶扩散}
\citet{Shay2004}以 $\B'=(1-d_e^2\nabla^2)\B$ 为推进变量，并用FFT恢复 $\B$。尽管密度在1--1.5之间变化，$d_e$ 在反演中作为空间常量；因此它是选定的惯性近似。默认网格间距 $0.1d_i$、质量比25给 $d_e/\Delta x=2$。所有方程加入四阶小尺度耗散，默认系数 $5.1\times10^{-5}$；高分辨算例同时把它改为 $5\times10^{-7}$。复现时要保存这个改动，不能只写“无碰撞、无电阻”。

\citet{Cassak2008}采用另一套设置：$204.8\times102.4d_i^2$、$\Delta x=0.05d_i$、质量比25，冷电子与绝热离子，四阶扩散约 $6.25\times10^{-6}d_i^3V_{\!A,\max}$。注意四阶扩散的系数与网格、单位和作用变量必须一起报告。两个都叫F3D的结果不一定采用相同的热力学或电子惯性。

\subsection{隐式Hall推进：稳定性收益与精度边界}
\citet{Arnold2008}把Hall部分作Crank--Nicolson推进，其他部分采用Runge--Kutta，利用局地 $3\times3$ whistler预条件及多重网格处理刚性。其GEM型测试为 $5\times2.5$ 区域、$256\times128$ 网格；报告时间步从 $0.0002$ 增至 $0.0082$，约41倍，但实际总时间节省约80\%，因为每个隐式步更贵。每个时间步需要约2--3个V循环是该算例的表现，不是所有参数的复杂度保证。

该预印本局地逆矩阵式(21)的 $(3,1)$ 元素印刷符号与原矩阵乘积不符；重新求逆时应检查乘积等于单位阵，不应机械复制。这里指出的是所获版本的印刷式，不能据此断言作者代码实现错误。

\citet{Toth2008}采用有限体积AMR与隐式Newton--Krylov相关推进。Hall电流需额外空间差分，面通量构造和幽灵单元深度随之改变；其实现需第三层磁场幽灵数据。该文还表明，直接给普通Lax--Friedrichs通量加上随 $1/\Delta x$ 增大的whistler速度，可能破坏一阶方案的一致性；重构与数值耗散必须一起分析。单波收敛中MC限制器和minmod的结果不同。重联测试中的大时间步加速需要与磁通曲线及结构一起核对，不能只比较程序是否稳定。

\subsection{五矩／十矩有限体积：电磁源耦合与压力保护}
\citet{Wang2015}演化全Maxwell双流体，通过局地隐式处理刚性源、有限体积波传播处理输运；正常高阶更新不保正时退回更耗散的Lax通量。该机制影响电子层，复现记录应包含回退触发位置与次数。局地源求解解除部分等离子频率限制，但传播仍有光速CFL。

其重点对照算例采用 $100\times50d_i^2$、质量比25、$\VA/c=1/15$，流体网格 $4095\times2047$，约8格解析 $d_e$；相应PIC网格更细。局地闭合设置 $k_{0e}=k_{0i}=1/d_e$，并非两物种都取各自惯性尺度。比较图中“等重联磁通”的流体与PIC快照发生在不同时间，不能解释成重联时史相同。压力非对角峰值也有显著差异；验收不应只保留最相似的磁力线图。

\citet{Hakim2008}的GEM例子采用十矩离子加充分松弛的五矩电子，而非两种粒子都保留无松弛十矩。其质量比1832的示例属于另一个Riemann问题，不能移植到GEM参数表。这是阅读方法论文时必须逐测试分辨的细节。

\subsection{热流闭合算法会改变物理结果}
\citet{Ng2017}用Fourier空间非局地闭合逼近动理学热流；\citet{Ng2020}用温度梯度形式并以超级时间步处理扩散刚性。两者的层厚、岛合并率尺寸标度和低杂漂移模增长率不同，说明闭合并非纯数值附件。

\citet{Closure2026}采用 $1152\times576\times288$ 三维网格、质量比100，并比较局地与梯度闭合；梯度限制器及闭合带来约3倍成本。其三维低杂波、kink与温度各向异性仍与Vlasov对照存在时序或空间位置差异，且缺少完整质量比／分辨率扫描。该预印本若干印刷公式的符号与对称化记号需要复核；本报告第一章用自行推导的一致方程说明模型，而不据疑似排版问题直接判断代码错误。

\subsection{三维扩展测量的两种诊断不能混合}
Shepherd--Cassak采用每个$z$切面的最大电流超过$3.3$标记重联区，并结合Hall电场阈值；Arencibia等改用离X点左右各$2d_i$处的平均$|B_y|>0.04$。后者的盒子为$102.4\times51.2\times256d_i$，网格$0.05\times0.05\times1d_i$，电子质量比25，导引场0或3。磁场种子系数0.005，$z$半长15，另以9、12、30检查；压力种子则检查机制是否随启动方式改变。\citep{Shepherd2012,L275}

Arencibia算例仍在所有演化方程加入四阶耗散：横向系数$1.6\times10^{-5}$，$z$向无／有导引场分别0.16／0.64；时间步分别0.02／0.01。各向异性网格和耗散必须逐方向报出。其v1式(17)把两侧高斯压力峰的两个平方项写在同一指数内，与文字描述不一致；复现压力启动须索取实际输入，而不能直接复制印刷式。

\subsection{陀螺流体的算法资料与证据缺口}
Grasso等的三场模型用双周期$[-\pi,\pi]^2$、$1024\times128$网格，$d_e=0.2$；平衡用$\operatorname{sech}^2x$的23个Fourier模。原文给出能量误差小于1\%的可信时间窗，却没有足够的时间推进、空间离散和网格扫描细节。因此能重建方程与初值，尚不能仅凭本篇保证逐图数值复现。Biancalani与Scott的GEM\emph{陀螺流体代码}另含温度矩与Landau闭合；这里GEM不是2001年的重联挑战。其最高$N_x=2048,N_y=512$，第三维点数未列，径向Dirichlet边界允许能量输入；不能把它当作全周期保守三场算例。\citep{E469,L184}

\subsection{Hall岛合并：网格改变的是方程中的有效耗散}
Bard与Dorelli的有限体积计算不含显式电阻或电子惯性，磁断线依赖数值耗散。其Fadeev盒为$4\pi\lambda\times2\pi\lambda$，$x$周期、$y$导体；$\lambda=6.04$--$25.46d_i$，部分扫描至$95.37d_i$，网格约$0.012$--$0.2d_i$。在$0.04$--$0.05d_i$附近层形显著改变，既体现尺度解析能力，也体现数值电阻改变。若要判定连续模型的物理解，应另加可控$\eta$或电子层机制，在固定物理参数下继续细化并测量欧姆残差。原文网格扫描本身不能替代这一步。\citep{L251}

\subsection{递归撕裂：先测局地层，再做标度}
Papini等使用$y$向Fourier、$x$向四阶紧致差分和三阶Runge--Kutta；磁势形式保持$\nabla\cdot\vct B=0$。初始力自由片半厚$a/L=S^{-1/3}$，$x\in[-20a,20a]$用特征投影自由流出，$y$周期且$L_y=100\pi a$；10个随机相位模的势扰动为$10^{-6}$。线性网格$1024\times128$、非线性$4096\times512$。

其局地诊断值得复用：从背景扣除后的$J_z$沿片剖面取半高全宽定义$2L_*$；在中间半段的多个位置拟合$B_y=B_*\tanh(x/a_*)$，用均值与标准差记录$B_*,a_*$；仅在$\Delta a_*/a_*\ll1$时把结构看作一张可拟合的片。随后用\emph{局地}$S_*=L_*B_*/(D_\eta\sqrt{\mu_0\rho_*})$而非原盒子的$S$做图。密度变化约1\%使该文可近似$\rho_*=1$，别的算例不能自动照用。\citep{L262}

\subsection{两电子群输入文件：用实际初始化复核参数}
2026年两电子十矩算例已保存三个公开输入文件于
\path{research/source_inputs/F635/}，来源为
\href{https://doi.org/10.5281/zenodo.21114117}{作者Zenodo数据记录}。
Gkeyll实际使用$480\times640$、$48d_e\times64d_e$、双周期、分维有限体积和CFL$2/3$，电磁源直接求解，每群有$k_0=0.1/d_e$的显式十矩松弛。两群初始密度之和恒为$n_0$，各自下限$0.03n_0$；这既避免真空，也是一项模型参数。

输入开头虽声明\texttt{plasmaBetae0=0.3}和\texttt{massRatio=25}，实际电子温度由$v_{\rm th}=0.111803c$设定，实际反转场为$0.5$、上游导引场$0.05$，所以$\beta_{e,\rm rec}\simeq0.1$。被推进及耦合的物种列表只有两个电子群；不能凭未参与演化的离子质量声明把它解释成有限质量离子计算。公开文件能补足论文叙述，但仍需检查代码版本、离散通量与输出处理，才可称完全复现。\citep{F635}

\subsection{装置几何中的谱元隐式计算}
MAST的HiFi计算使用每单元四阶Jacobi基函数、非均匀谱元网格和隐式Crank--Nicolson。常用$180\times360$单元，对应$720\times1440$有效分辨率，片附近最小间距$0.23$ mm；最低超电阻例用$360\times540$单元。Hall时间步上限为$0.01\tau_0$，$\tau_0=0.29\,\mu$s，隐式求解仍因精度要求限制步长。

环形域$R=0.2$--$2$ m、$Z=\pm2.2$ m，导体壁条件与垂直背景磁通同时处理；平面算例用滑移壁，环形算例因法向磁场改用无滑移。初值从磁绳已脱离线圈后开始，不模拟气体击穿及绳形成。最低超电阻的双分辨率峰率差为0.2\%，但比较时间截止次级岛出现前；这一检查不自动认证后期湍动统计。\citep{Stanier2013}

\section{全隐式Hall求解：从非线性残差到可扩展的预条件}
\subsection{隐式步究竟求解什么}
以半离散方程$\dot{\vct U}=\mathcal L(\vct U)$为例，Crank--Nicolson要求解
\begin{equation}
 \vct G(\vct U^{n+1})=\vct U^{n+1}-\vct U^n
 -\frac{\Delta t}{2}[\mathcal L(\vct U^{n+1})+\mathcal L(\vct U^n)]=0.
\end{equation}
BDF2则把时间导数换成
$(3\vct U^{n+1}-4\vct U^n+\vct U^{n-1})/(2\Delta t)$。
Newton每次求$\mathsf J\delta\vct U=-\vct G$，再更新$\vct U$。直接存储三维网格的巨大Jacobian很昂贵；Jacobian-free方法只在Krylov迭代需要时计算
\begin{equation}
 \mathsf J\vct z\simeq
 \frac{\vct G(\vct U+\epsilon\vct z)-\vct G(\vct U)}{\epsilon}.
\end{equation}
$\epsilon$要兼顾截断与浮点消去误差，并根据变量尺度缩放。这里仍要完整计算物理残差，不能因不存矩阵就漏掉耦合项。\citep{L296}

\subsection{块分解：预条件近似不等于删去物理}
把密度、压力、矢势等合为$\vct y$，离子速度为$\vct v$，简写线性系统为
\begin{equation}
 \begin{pmatrix}\mathsf M&\Delta t\mathsf U\\
 \mathsf L&\mathsf D_v\end{pmatrix}
 \begin{pmatrix}\delta\vct y\\\delta\vct v\end{pmatrix}
 =-\begin{pmatrix}\vct G_y\\\vct G_v\end{pmatrix}.
\end{equation}
若精确消去$\delta\vct y$，速度块为
$\mathsf S=\mathsf D_v-\Delta t\mathsf L\mathsf M^{-1}\mathsf U$。
实际用可便宜反演的$\widetilde{\mathsf S}$和近似$\mathsf M^{-1}$，依次得到临时$\delta\vct y$、速度修正和回代。这是加速寻找原系统解的方法；只有Newton最终残差也达到阈值时，才可把预条件中的近似与物理截断分开。

Chacón的做法依据离子／电子时间尺度分开块结构：先处理电子Hall电磁块，再处理离子流动耦合；小体流等近似用于构造Schur预条件。外层采用\emph{flexible} GMRES，因为内层GMRES的停止次数会变，预条件不再是每次完全相同的线性算子。原文Newton相对容差$10^{-3}$、绝对容差$10^{-15}\sqrt N$，Krylov误差随Newton残差调整，以免早期线性问题解得过精。

\subsection{为什么额外引入电流未知量}
对矢势直接代入$\vct j=\nabla\times\nabla\times\vct A$，电子黏性会产生四阶空间算子。把$\delta\vct A$和$\delta\vct j$各三分量同时求解，把四阶问题改为耦合的二阶块，便于多重网格平滑。该文采用V(3,3)循环、阻尼系数0.7的块Jacobi；二维四颜色九点、三维八颜色二十七点模板。限制为守恒聚合，延拓为一阶，粗网格算子重新离散。电磁内层GMRES最多10步或相对残差降到0.1；其他较易子块通常各一个V循环。

恒等式$\nabla\times\nabla\times\vct A=\nabla(\nabla\cdot\vct A)-\nabla^2\vct A$说明，只在Coulomb规范下才可普遍删掉第一项。这里选Weyl规范，作者仅在\emph{Jacobi局地对角近似}里用$-\nabla^2$，完整残差仍保留curl--curl。把这个近似误移到实际方程会解成另一个问题。

\subsection{性能数字和验证范围}
波传播用Crank--Nicolson，GEM用BDF2和二阶有限体积；$\B$由兼容离散旋度得到以控制磁散度。GEM物理比较为$256^2$拉伸网格、$\Delta t=0.05$，与HiFi多数时段能量曲线相近，末期约1\%差。所测试参数没有形成中心岛，只是这一组耗散和边界下的观察。

\begin{center}\small
\begin{tabular}{rrrr}\toprule
网格 & MPI任务 & $\Delta t/\Delta t_{\rm explicit}$ & 实际运行加速\\\midrule
$32^2$ &1&1&0.04\\
$256^2$&64&33&0.63\\
$512^2$&256&130&3\\
$4096^2$&16384&8400&75\\\bottomrule
\end{tabular}\end{center}
每个MPI任务固定$32^2$单元，因此这是\emph{弱扩展}。最高8400是允许时间步之比，不是运行速度之比；小网格隐式反而慢。性能扫描的电子黏性从$10^{-4}$随网格减到$6\times10^{-9}$，所以它不是固定方程的物理收敛试验；100步、$\Delta t=0.01$也仅覆盖$t\leq1$，早于所示GEM的强非线性阶段约$t=20$。不能据此保证磁岛级联全过程也同样高效。

接收稿附录的平滑分析限制在均匀周期背景、零电子流、$d_e=0$和刚性且有足够网格黏性的情况；中心差分的某些零符号模式也需要单独讨论。它提供预条件设计依据，没有证明任意非均匀、无黏性和三维问题都具有相同收敛速度。与所有隐式算法一样，稳定跨过快速波周期不代表这些波的相位或电子层动力学仍被准确解析。

\section{递归Hall谱代码：需补齐的复现信息}
Shi等的另一套计算是双周期伪谱、显式三阶Runge--Kutta、$8\times0.8$盒和$4096\times2048$网格；它与Papini等混合紧致差分／Fourier方案不同。均匀电阻$10^{-5}$、初片半厚0.02，七组$d_i$从0至0.01，随机相位1--128模在各组完全相同。论文未给种子幅度、完整去混叠／滤波和步长系数，不宜声称仅凭已列参数即可逐图复现。

其谱横轴采用循环频率$N/L_x$，不是通常的角波数$2\pi N/L_x$；讨论$kd_i=1$时需先换约定。磁通率取固定子区域内主X点的磁势随时间线性拟合，没有展示每个O点参考势项；跨拓扑变化时应另外检验净重联磁通。\citep{L263}

\section{受迫与强导引场算例：需要逐篇保存的实际输入}
\subsection{Taylor问题的半隐式计算}
两篇Taylor研究均用$256\times128$网格、$\Delta t=10^{-3}$，$x$方向在电流层附近加密，空间与时间二阶，使用Harned--Miki\'c型半隐式处理。$x=\pm1$导体墙、$y$周期长度8，驱动上升时间1；常用超电阻$\nu=2.5\times10^{-9}$。任意导引场篇另取黏性和热扩散$\mu=\kappa=10^{-7}$。这些量按各文的系统尺度归一化，不能与F3D按$d_i$归一化的耗散系数直接比较。\citep{L057,L058}

两文给出了耗散与驱动扫描，但没有独立的网格、步长收敛系列。复算应先测X点电阻和超电阻的贡献，再在固定物理系数下细化；若同时降低超电阻，会把模型变化与离散误差混在一起。墙位移、墙速度与磁通边界必须彼此一致，具体关系已在第一章推导。

\subsection{F3D受迫算例：电子惯性不替代X点耗散账目}
2005年二维算例用$102.4d_i$方盒、$512\times1024$网格，$\Delta x=0.2d_i,\Delta y=0.1d_i$、质量比25；质量比100对照用$1024\times2048$。空间四阶差分，时间二阶trapezoidal leapfrog；密度和动量加四阶耗散$5.1\times10^{-5}$。2008年三维改为$25.6d_i$周期立方体、$128\times256\times128$，仍只有约两格跨法向$d_e$。\citep{L076,F717}

模型用常系数$\B^*=(1-(m_e/m_i)\nabla^2)\B$，并非一般变密度电子惯性。准稳态X点仍依赖数值扩散；所以“没有显式电阻”不能写成“X点非理想项全由惯性支撑”。三维没有给出充分分辨率扫描，第三维较粗也限制不稳定性研究。

复现前还需处理原文输入疑点。2008年所获版本的驱动时间函数印作$\tan$，与稳定渐开驱动的叙述及2005年的$\tanh$形式矛盾；$\tau_f$也有前后不同写法。这些应向输入文件追溯，不能用猜测的公式宣称逐图复现。式(17a)、(18a)的出流标度疑点则可由连续性与动量平衡独立发现。

\subsection{PIXIE2D与VPIC：局地对照和尺寸扫描用了不同参数}
\citet{L217}使用Newton--Krylov、物理多重网格预条件和向X点集中的网格；空间二阶中心差分，对流项用QUICK，时间可选Crank--Nicolson或BDF2。论文未逐算例列出格点数、最小网格和时间步，因此只能登记其方法与收敛陈述，不能补造完整输入。

\begin{center}\small
\begin{tabular}{lrrr}\toprule
VPIC比较 & 无快速色散波局地例 & 有快速色散波例 & 无快速色散波尺寸例\\\midrule
$m_i/m_e$ &15.2&380.3&15.2\\
$B_g/B_{\rm rec}$ &17.4&12&17.4\\
$v_{the}/c$ &0.0625&0.308&0.25\\
$\omega_{pe}/\Omega_{ce,x}$ &10&10&2.5\\
$T_i/T_e$ &0.25&0.1&0.25\\\bottomrule
\end{tabular}\end{center}
尺寸例的Alfvén归一化还考虑相对论修正。长文图14与15的流体黏性分别为$10^{-4}$与$10^{-5}$；有快速色散波的PIC例每格16000粒子。这些细节说明，不能把所有图拼成同一质量比、同一温度和同一耗散的单变量扫描。\citep{L216,L217}

\section{陀螺流体和非局地热流：比网格大小更容易漏掉的设置}
\subsection{四场双重加速的伪谱计算}
\citet{F672}把$0.1\operatorname{sech}^2x$平衡截成23个Fourier模，采用双周期边界、三阶Adams--Bashforth和伪谱求导，最高$4096\times512$，并在微尺度滤波。非线性系列$L_y=4\pi$、$k_y=0.5$、$\Delta'\simeq59.9$，$d_e=0.2$至0.025，同时保持$d_i=10d_e,\rho_s=2d_e,r_i=4d_e$。

所报告的一次$x$分辨率加倍，峰值率基本不变、峰值电子流改变不到6\%；它支持该例宏观率的稳健性，不等于所有细丝、所有参数均收敛。微层继续变细到不能解析时应停止判读。原文线性图2在正文写$\rho_s=0.1$，图注写0.2；在获得输入前，该图不能作为精确参数复现的验收样本。

\subsection{陀螺动理学对照：边界、种子和模过滤都要登记}
\citet{E475}的流体实现使用$1024\times128$、双周期伪谱；EUTERPE使用B样条有限元与四阶Runge--Kutta，$x$为Dirichlet而$y$周期，最高约$3\times10^7$标记粒子、最小时间步0.125，$d_e$至少16格。时间归一化须按该文定义换算，不能与前一文的步长直接相比。

流体用确定种子，PIC由噪声开始。边界差别通过局地线性本征模检查受到约束，但非线性比较仍应记录。PIC非线性保留$m=0,\pm1$，流体保留更高谐波；作者认为小岛区高谐波很弱，这不能外推到大岛分裂或双重加速。报告的总能量误差小于2.5\%也应换算成释放磁能比例后，才适合评估加热精度。

\subsection{非局地闭合：刚性与三维模空间必须分开验证}
\citet{L244}用三阶Adams--Bashforth推进，沿周期方向并行。非线性加入热流后，步长可需缩小约100倍；这是所测参数的相对成本，原文没有列出绝对步长和完整网格数量。非局地分母用磁场方向二阶矩的空间平均，分子仍含局地$\nabla_\parallel$，不能改成普通常系数Laplacian而称实现了同一闭合。

线性半盒$(L_x,L_y,L_z)=(46,24\pi,96\pi)$，非线性用$k_y=1/12,k_z=1/48$及$(L_y,L_z)=(12\pi,48\pi)$。单螺旋性允许用$y^*$二维坐标降低成本，却也删去了不同螺旋性的耦合。验证应分别检查闭合算子的线性阻尼、单模增长、能量损失，以及加入多模后的三维演化；任何一个检查不能替代其余检查。

\section{五矩双流体的两条可实施路线：DG与全隐式有限体积}
\subsection{DG：一个单元内保留多个自由度}
\citet{L156}将两物种Euler方程与完整Maxwell系统写成
$\partial_tU+\nabla\cdot F(U)=S(U)$，每单元用低阶多项式展开$U_h=\sum_jU_j(t)\varphi_j$。
乘试验函数、单元积分并分部积分后，有
\begin{equation}
 M_{ij}\dot U_j=\int_K\nabla\varphi_i\cdot F(U_h)\dd V
 -\int_{\partial K}\varphi_i\widehat F\dd S+\int_K\varphi_iS(U_h)\dd V.
\end{equation}
正交基使质量矩阵$M$对角；单元间通过数值通量$\widehat F$交换信息。
该文一次基用于二阶空间精度，二次基用于三阶精度；二维基分别包含$\{1,x,y\}$及增加的$xy,x^2-1/3,y^2-1/3$。
边积分分别用2或3个Gaussian点、体积分4或9个点；局部Lax通量按物种声速和流速选耗散，电磁部分含光速。

限制器也是算例定义的一部分。该文使用特征空间minmod/TVB；一旦线性系数被限幅，高阶系数清零，因此“采用三阶DG”不意味着每个单元始终三阶。
时间采用SSPRK3统一推进通量与刚性源项。对纯振荡$\dot q=i\omega q$，二阶显式RK的放大因子模大于1，不能以流体平流CFL代替源项稳定性分析。
文中典型要求$\Delta t<\omega_{pe}^{-1}$，还受光波CFL约束。所读预印本第12页三阶CFL的$c/\Delta x$量纲错误；应具有$\Delta x/c$量纲，具体安全系数随基阶和空间维数而变。

GEM部分使用$m_i/m_e=25$、$c/\VA=10$、$T_e/T_i=0.2$，比较$128\times64$与$512\times256$网格，时间步分别$0.002$和$0.0005\Omega_{ci}^{-1}$。
粗格二阶持续出岛而三阶没有同样的岛；细格两阶通量接近。更强限制/耗散也可能消岛，说明不能给数值耗散与磁岛数建立简单单调关系。
晚期局地动量对浮点精度有约10\%差别，且磁岛偏移方向可受舍入误差影响。
这要求在物理结论旁报告网格、阶数、限制器、时间步和误差敏感性，而不只列网格数。

\subsection{隐式有限体积：稳定步长、矩阵内存与求解成本}
\citet{L243}采用18变量五矩双流体--Maxwell--散度清理系统。
二阶空间重构使用最小二乘梯度及Venkatakrishnan限制器，物种通量用AUSM+up，电磁通量用改进Rusanov；隐式BDF2配Newton迭代，数值差分组装Jacobian，再以GMRES/Jacobi求解。
它实际组装矩阵，不能笼统归入无矩阵Newton--Krylov算法。

GPU实现以一个线程处理一个单元，计算通量、源项与变量扰动后的Jacobian块；CSR矩阵装配移到设备侧，避免反复CPU传输。
边界处理等仍有主机工作。波算例检查二阶误差；GEM采用双电流片全周期边界、$256^2$与$512^2$单元、质量比25，至$t\Omega_{ci}=25$共10000步。
后者$\Delta t\Omega_{ci}=0.0025$，P100完整一次计算约49小时16分；没有因“全隐式”就无限加大物理时间步。

报告的14.23倍加速相对同机\emph{单个CPU核心}；若换更快的K40主机CPU作基准，约为8.15倍。
小问题GPU占用低，某些PARALUTION配置比CPU慢；大问题又受到完整耦合矩阵显存限制。这些是当时PETSc 3.9、PARALUTION 1.1与特定硬件的结果，不能用于预测今日任意平台。
GPU/CPU磁通吻合验证实现一致性，尚不验证电子压力闭合或局地非理想电场的物理准确性。

\section{线性、非线性与湍流算例的实际分辨要求}
\begin{longtable}{P{29mm}P{64mm}P{55mm}}
\toprule 文献与路线 & 已核实的计算细节 & 必须保留的限制\\\midrule\endhead
\citet{L224}：斜向惯性撕裂 & HiFi谱元/BDF2；旋转二维平面模拟不同倾角；$m_i/m_e=100$，$d_e$约8--11格，导场比1--10，超电阻扫描 & 各角度独立二维线性计算，不等于一次三维非线性竞争\\
\citet{L268}：电阻Hall本征值 & 不可压Harris平衡；SciPy四阶配点，横向区间约$[-15,15]$，同时扫描两方向波数 & 是本征边值问题，未列统一固定网格/容差；增长与传播频率应分别核对\\
\citet{L191,L214,F697}：爆发阶段 & 早期$x$约$10^4$差分点、$y$约200模；后续二维谱/RK4达到$4096^2$--$8192^2$ & 增加横向模数后看到次级岛，旧的平滑加速结论受分辨率影响\\
\citet{E471,E472}：回旋流体 & 通常$1024\times128$，部分非线性提高到256横向模；Padé回旋算子 & 压力/导心密度能不同于热化；总能误差约百分数不能替代热预算\\
\citet{E474}：有限$\beta_e$ & 谱空间、三阶Adams--Bashforth、短尺度滤波，最高$2880^2$，能量分析限于损失小于1\%时段 & 滤波后段不宜推断物理耗散；移动X点影响增长率曲线\\
\citet{F644}：KH卷涡 & 2.5维Hall，$4096\times8192$，盒约$400\times402d_i$，RK3，滤除$kd_i\gtrsim3$ & 无显式电阻、惯性及张量破冻结项；拓扑改变受滤波/离散影响\\
\citet{L294}：非均匀片传播 & F3D，$102.4\times51.2\times256d_i$，网格间距$(0.05,0.05,1)d_i$，质量比25 & 前沿判据$B_y=0.04$与非均匀片几何是拟合的一部分\\
\bottomrule
\end{longtable}

\subsection{嵌入Hall区时，过渡函数本身会进入方程}
\citet{L035}将局部Hall计算嵌入外部MHD，比较显示磁通演化接近全域Hall。
但若直接把非理想电场写成$\alpha(\vct x)\E_H$，感应方程中出现
\begin{equation}
 -\nabla\times(\alpha\E_H)=-\alpha\nabla\times\E_H-\nabla\alpha\times\E_H.
\end{equation}
第二项是过渡带带来的额外作用。论文给出的平滑erf混合不等于已证明守恒的多尺度耦合算法。
未来复现应同时改变Hall区域大小和过渡带宽度，检查过渡区能量及电流，而不只比较一个全局通量曲线。

\subsection{三种“率”应在数据文件里就分开}
保存$E_z/(B_mV_{\!A,m})$、$\dd\ln\psi/\dd t$和$\delta/L$三个独立量。
\citet{L264}的重联频率由磁通演化时标获得；\citet{F644}的0.1--0.3主要是电流片纵横比估计。
只有满足稳态、不可压、明确的出流速度和上游归一化时，纵横比才与电场率相联系。
在合并、湍流和非稳态小片中，这些量即使数值接近也不可互换。

\section{续筛新增的数值证据：先比较算法，再比较现象}
\label{sec:remaining-numerics}
本节列出此次全文核读直接影响判断的计算条件。所有物理推导在第一章；这里讨论实际求解与验证，不能把理论写出的方程等同于代码中实际保留的项。

\subsection{全局尺寸扫描应固定什么}
Huang等\citep{F600}用$L_x=25.6\ell,L_z=12.8\ell,\delta=0.5\ell$的Harris片，改变$\ell/d_i$，因此保持的是宏观形状和$\delta/L$，不是固定微观片厚后单独加长盒子。PIC扫描到$\ell/d_i=30$，Hall MHD到100。PIC使用质量比25、20格/$d_i$及每种每格平均250粒子；Hall流体在中心区用$\Delta x=0.1d_i,\Delta z=0.02d_i$，超电阻和黏性采用固定离子尺度归一值。

比较时间也须一致：相同$t\Omega_{ci}$不等于相同全局演化阶段。该文比较相同重联磁通区间并以初始$B_0V_A$归一。但实际核读v1的正文平均区间写$0.05$--$0.3$，图1阴影写约$0.5$--$3$；两者相差十倍。在原始脚本或作者澄清前，可引用已展示的尺寸趋势，不能自称精确复算其平均率。另有一个PIC尺寸的重复运行相差约12\%，应同时记录随机扰动敏感性。

磁岛合并的尺寸扫描\citep{L241}同样改变$\lambda/d_i$并保持形状，但系统由宏观岛供应磁通，与长Harris片并不相同。其率使用$B_m=0.353B_0$与$V_{Am}$归一，平均窗口为$1.5\tau_A$。直接把其约0.6的峰值与其他文章的0.1放在一起，会混淆磁场基准、局地/全局和峰值/平均。

\subsection{非对称层中，改变耗散项可以改变整个宏观演化}
Aunai等\citep{F625}用相同混合模型更换网格控制、电阻和超电阻处理。离子由Boris推进，电子无质量且等温，电磁场二阶差分和预测校正。非对称两侧$|B|=0.5,1.5$，密度$1,1/3$，导引场1；$T_i/T_e=5$。其结果不能简单写成“Hall打开了快重联”：
\begin{enumerate}
 \item 网格控制时，非对称电子层反复伸长并数值撕裂；细化网格改变起岛时刻、层长和平均率。对称算例较不敏感，不能代表所有几何。
 \item 均匀电阻可保持短层，但同时扩散更大尺度的磁能；超电阻更集中在细尺度，其宏观率也不同。
 \item 周期域太短时，返回的流和场可加宽电流层，使磁岛不再形成，或者只出现一次后便维持短X线。看起来稳定的解可能由边界支持。
\end{enumerate}
这是对“任意小尺度断线机制都等效”的反例，不是对某个超电阻闭合已等价于真实动力学的证明。该文部分加密组同时改变了盒长，故严格复现还需分别改变网格与域长，并只比较边界返回影响之前的时间窗。

\subsection{磁岛与分辨率：两类对照必须分开}
Garcia Morillo与Alexakis的正式版本\citep{F788}研究的是二维不可压\emph{电阻MHD}，并没有Hall项。双周期$2\pi$方域、Orszag--Tang初值、$\nu=\eta$、初始磁/动能$0.5/0.125$，另加$k<16$、能量约0.25\%的扰动；GHOST伪谱、四阶Runge--Kutta、$2/3$去混叠，$N=1024$到32768。欠分辨时出现磁岛和近乎平台的率；足够分辨的相应算例中没有这些岛，率更接近Sweet--Parker标度。最高已分辨局地$S_L$约$5.4\times10^5$。作者还讨论了电流片寿命、扰动衰减和被排走的时间，正式版并未证明“所有高$S$磁岛都是假的”。

该文用电流谱在截断波数处相对峰值充分衰减作实用检查，但其经验谱阈值不能直接移植为有限体积或DG的统一收敛判据。本轮已下载并核读10页期刊版本，原8页预印本保留，独立论文仍只算一篇。

Akutagawa等\citep{F606}则比较PIC与电阻MHD。主要PIC质量比1、2、4、8、16、32，空间步长$0.1d_e$，$100d_i\times50d_i$，部分加倍域；主文低质量比易出现次级岛，但这些岛并未像MHD例那样显著提高全局能量转换。附录显示一些粗网格磁岛在加密后消失。真实质量比结论仍是外推：附录虽给质量比100，使用的是$0.5d_e$粗网格，也不等于完成质量比1836验证。

两篇文章分别涉及不同方程、噪声和形成过程，不能合并成一个否定磁岛机制的普遍结论。更有用的共同要求是同时查看磁通拓扑、场谱、网格/时间步、扰动、边界返回、守恒和非理想项预算，而不是仅比一张电流密度彩色图。

\subsection{回旋流体的制造解、超黏性和实际算例是三种检查}
GREENY v1\citep{E473}给出极化和Ampère反演的制造解，展示SOR约$h^2$、其FFT或PCG/Teague组合约$h^4$误差；这里的阶数包含其具体离散，不表示所有谱方法只有四阶。算法使用Arakawa括号和Karniadakis多步，极化求解与磁势求解可选择不同近似。

近边界重联算例中，超黏性太强会抹平初始扰动，太弱会出现格点颗粒结构；部分主峰重合而前期岛形成和起始时间仍变。必须区分：制造解是否正确、守恒误差是否小、特定非线性现象是否随网格及耗散一起收敛。不能只展示其中一项便宣布物理已收敛。后续Full-F研究\citep{L314}的$1024^2$与$32\rho_s$域及迎风/超扩散比较提供了进一步信息，但文中若干参数标注矛盾仍需保留，见第\ref{sec:remaining-audit}节。

\subsection{近期大尺度混合湍流：电子尺度分辨尚需特别核验}
F655\citep{F655}采用$100\pi d_i$方域、$3072^2$空间格、$51^3$离子速度格，质量比100，等温惯性电子，初始注入尺度约$50d_i$。它确实扩大了相对注入尺度，但由公开数值可直接算
\begin{equation*}
 \Delta x={100\pi\over3072}d_i\simeq0.1023d_i,
 \qquad d_e=0.1d_i,\qquad\Delta x/d_e\simeq1.023.
\end{equation*}
电子惯性长度约一格，且使用格点尺度滤波，正文没有独立加密/滤波收敛组。宽度数个$d_i$的涡链可见，并不自动保证断线电子层已收敛。报告因此保留其跨尺度生成机制，同时降低电子层功率峰值和发生率外推的证据等级。

\subsection{RFP计算应把壁面通量写入验证}
NIMROD扩展MHD松弛研究\citep{F798}用$r$--$z$面$120\times64$个双五次谱元、方位$m=0$--5、轴向周期，比较有无Hall及离子回旋黏性；$S=20000$、$d_i/a\simeq0.17$，电子质量放大5倍。它没有求完整热预算，温度方程忽略电阻与黏性加热。

新生成的总动量主要来自壁面的黏性与回旋黏性耦合，再由Maxwell应力向内输运。作者还量化了省略平衡pinch流导致的额外动量损失；特定事件中相对于$r/a=0.7$处Maxwell通量约10\%到22.5\%。因此核心速度变化方向的比较有价值，但不能把总动量增长当作已解释了真实实验的无外源自发加速。

